\section{Optional Cavalry Momentum Rule}

\begin{rulebox}{Purpose}
The cavalry fighting at Marston Moor was capable of carrying an entire wing
away with it. This rule lets Heavy Horse gain ground after a successful combat,
without adding a second attack or a full pursuit system.
\end{rulebox}

\begin{multicols}{2}
\subsection{1. Cavalry recoil}

When a combat involving at least one attacking ordered Heavy Horse unit inflicts
a Dd or Dx result on one or more ordered defending Horse units, apply the
disruption and then retreat each newly disrupted defending Horse unit one hex.
Heavy and Light Horse both recoil. A defending Horse unit that was already
disrupted is eliminated normally and does not retreat.

The defending player chooses each retreat hex. It must be an adjacent, on-map
hex that contains no unit and that the retreating Horse could legally enter by
normal movement. The unit may not cross a prohibited hexside. Movement-point
costs and enemy Zones of Control are ignored. If possible, the retreat hex may
not be adjacent to any unit that participated in the attack. If no legal hex is
available, the unit remains where it is, disrupted.

Any leaders stacked with a retreating unit may accompany it. Artillery pieces
never retreat under this rule.

\subsection{2. Heavy Horse advance}

After all results of a combat, including cavalry recoil, have been applied, an
ordered Heavy Horse unit that participated in the attack may advance into an
adjacent hex vacated by a defending unit in that combat. If more than one
defending hex was vacated, one eligible Heavy Horse unit may advance into each
vacated hex. No unit may advance more than once per Combat Phase.

An advance is optional, costs no movement points, and may enter an enemy Zone
of Control. The advancing unit must obey normal terrain, hexside, and stacking
restrictions. Any leaders stacked with it may accompany the advance.

An advancing unit cannot attack again or add its strength to an unresolved
combat. Adjacencies and Zones of Control created by an advance do not require
the players to declare additional combats during the current Combat Phase.

\subsection{3. Charges}

A charging Heavy Horse unit may advance if it remains ordered after the combat
result. Make the advance before applying the automatic disruption required by
\ruleref{10.3}; the unit is then disrupted in its new hex. The directional
Charge marker may be placed on a charging unit with its arrow pointing toward
the defending hex.

\begin{rulebox}{Sequence}
\begin{enumerate}[leftmargin=1.3em,itemsep=.12em,topsep=.1em]
  \item Apply the combat result.
  \item Recoil newly disrupted defending Horse.
  \item Advance eligible Heavy Horse into vacated hexes.
  \item Disrupt ordered Heavy Horse that charged.
\end{enumerate}
\end{rulebox}

\section{Historical Note}

The historical record supports a more fluid treatment of cavalry. Goring's
cavalry broke Fairfax's wing, with some squadrons pursuing or plundering while
others turned against the Allied infantry. On the other flank, Cromwell and
Leslie eventually drove off the Royalist horse; Cromwell then kept enough of his
command together to cross the battlefield and attack the Royalist centre and
Goring's scattered troops. The cavalry action produced sudden changes of
position, even though the fight against Rupert's horse was itself
hard-fought.\footnote{See
\href{https://historicengland.org.uk/listing/the-list/list-entry/1000020?section=official-list-entry}{Historic England, ``Battle of Marston Moor 1644''}
and
\href{https://www.battlefieldstrust.com/resource-centre/battlepageview.asp?pageid=698}{Battlefields Trust, ``The Battle of Marston Moor.''}}

The recoil gives the winning side ground without treating every disruption as
destruction. Limiting the follow-up to Heavy Horse preserves the shock role of
that arm. The prohibition on another attack keeps the result to a local
breakthrough rather than allowing cavalry to roll up several units in one
Combat Phase.

\begin{rulebox}{Rules References}
Combat declaration: \ruleref{9.0}. Heavy Horse charges: \ruleref{10.2--10.3}.
Combat results and disruption: \ruleref{11.0--11.8}.
\end{rulebox}
\end{multicols}
